\Needspace{4\baselineskip}
\begin{Example}
\label{ex:sobolev-first-second-norms}
For $1\leq p<\infty$ we have\[
\|u\|_{1,p}=\left\{\int_{\Omega}\left(|u|^p+\sum_{i=1}^m\left|\dfrac{\partial u}{\partial x_i}\right|^p\right)\,\mathrm{d}\bm{x}\right\}^{1/p},
\]
\[
\|u\|_{2,p}=\left\{\int_{\Omega}\left(|u|^p+\sum_{i=1}^m\left|\dfrac{\partial u}{\partial x_i}\right|^p+\sum_{i,j=1}^m\left|\dfrac{\partial^2u}{\partial x_i\partial x_j}\right|^p\right)\,\mathrm{d}\bm{x}\right\}^{1/p}.
\]
\end{Example}

\begin{Example}
\label{ex:heaviside-absolute-value-sobolev-membership}
The Heaviside function $H(x)$ is in $L^p(-1,1)$ for any $p\in[1,\infty]$. However, it is not in $W^{1,p}(-1,1)$ since its distributional derivative, the Dirac delta, cannot be represented by an $L^p(-1,1)$ function. Similarly, the absolute value function $f(x)=|x|$ is in $W^{1,p}(-1,1)$, since its distributional derivative can be represented by the $L^p(-1,1)$ function\[f'(x)=\begin{cases}-1&x<0\\1&x>0\end{cases}\]However, $f$ is not in $W^{2,p}(-1,1)$ since its second distributional derivative $(2\delta)$ cannot be represented by a function in $L^p$.
\end{Example}

\begin{Example}
\label{ex:radial-power-sobolev-membership}
Consider the function\[f(\bm{x}):=|\bm{x}|^{\alpha}=\left(\sum_{i=1}^m x_i^2\right)^{\alpha/2}\]on $B=B_1(\bm{0})\subset\mathbb{R}^m$, the ball of radius one in $\mathbb{R}^m$. One can easily check using spherical coordinates that if $\alpha+m>0$ this function is integrable. Thus, if $p\alpha>-m$ we have $f\in L^p(B)$.

Formally taking partial derivatives we get\[\dfrac{\partial f}{\partial x_i}(\bm{x})=\alpha x_i\left(\sum_{i=1}^m x_i^2\right)^{\alpha/2-1}=\alpha|\bm{x}|^{(\alpha-2)/2}x_i\]\[
\dfrac{\partial^2f}{\partial x_i\partial x_j}(\bm{x})=\alpha(\alpha-2)x_ix_j\left(\sum_{i=1}^m x_i^2\right)^{\alpha/2-2}+\alpha\left(\sum_{i=1}^m x_i^2\right)^{\alpha/2-1}\delta_{ij}
\]where $\delta_{ij}$ is the Kronecker delta.

We can use standard techniques of improper integrals to show that if the formal derivative has an integrable singularity, then that function is a representative of the distributional derivative. Using this and calculating the norms we get (after some work)\begin{itemize}
\item $f\in W^{1,p}(B)$ if and only if $p(\alpha-1)>-m$ (i.e., $\alpha>1-m/p$).
\item $f\in W^{2,p}(B)$ if and only if $p(\alpha-2)>-m$ (i.e., $\alpha>2-m/p$).
\end{itemize}Similar calculations show that $|\bm{x}|^{\alpha}\in W^{k,p}(B)$ if and only if $\alpha>k-m/p$.
\end{Example}
